\itemhead{p.84 练习 · 第2题（教材p.84）}{XJ-B1-C03-PR-P084-Y1402}
\label{XJ-B1-C03-PR-P084-Y1402}
{\footnotesize 湘教版教材p.84 / PDF第91页 / p.84 练习 / 第2题}
\field{原题}{函数 $f(x)=ax^2+bx+c$ 在什么条件下是偶函数？在什么条件下是奇函数？在什么条件下既是奇函数又是偶函数？}
\field{参考答案}{不预设 $a\ne0$ 时：偶函数当且仅当 $b=0$；奇函数当且仅当 $a=c=0$；既奇又偶当且仅当 $a=b=c=0$。若额外限定为二次函数（$a\ne0$），则不可能是奇函数。}
\field{核对说明}{版本插题变化：现教材增加多项式奇偶条件，答案册PDF23同组后面乘积、奇延拓题号与现题不同。}
{\small \field{关联知识}{F08 奇偶性、对称与函数性质应用（教材帮p.113、114、115 / PDF122、123、124）}
\field{方法与易错点}{检查定义域，再代入 $-x$；求另一侧解析式时同时写出 $-x$ 所在区间。 易错：只看公式不看定义域；用一对点证明奇偶；负半轴代换遗漏最外层负号。}
\field{答案依据}{答案册相关题组：PDF第23页（册内页22）；按当前题库原题逐项重算，不按旧版题号机械对套。}}
